\par\bigskip
\Needspace{10\baselineskip}
\hypertarget{native-373-lower}{}\label{native-373-lower}
\begin{center}\large PALLAS.\end{center}
\begin{center}
\parbox{.73\linewidth}{%
\centering\fontsize{8.2}{10.2}\selectfont
\hrule height .3pt\relax\vspace{2.5pt}
Monatliche Correspondenz. Band XXVIII. Seite 197..198. 1813 August.\par
\vspace{2.5pt}\hrule height .3pt\relax
}
\end{center}

\begin{flushright}
Göttingen, den 31. Aug. 1813.
\end{flushright}

\par\indent --- --- --- Ich übersende Ihnen hier meine in diesem Jahre erhaltenen \emph{Pallas}-Beobachtungen:

\begin{center}\fontsize{10}{12}\selectfont\setlength{\tabcolsep}{5pt}\renewcommand{\arraystretch}{1.15}\begin{tabular}{l|c|c|c}
1813 & \shortstack{Mittlere Zeit\\in Göttingen} & AR. $\PallasSymbol$ & Decl.\\\hline
Juni 28 & $11^{\mathrm u}\;37^{\mathrm m}\;2^{\mathrm s}$ & $327^\circ\;2'\;13''3$ & $+14^\circ\;5'\;43''4$\\
Aug. 16 & $10\;36\;23$ & $319\;50\;54.7$ & $.\;.\;.\;.\;.$\\
19 & $9\;31\;0$ & $319\;17\;10.6$ & $10\;0\;8.4$\\
26 & $9\;23\;7$ & $317\;59\;52.0$ & $8\;44\;35.3$\\
\end{tabular}\end{center}

\par\indent Herrn \textsc{Nicolai}'s Vergleichung mit den letzten Elementen gibt folgende Resultate:\par
\clearpage
\thispagestyle{empty}
\noindent\SourceLeaf{374}%
\makebox[12mm][l]{}%
\parbox[t]{\dimexpr\linewidth-24mm\relax}{\centering\footnotesize\bfseries PALLAS.\par}%
\makebox[12mm][r]{{\large\bfseries 369}}\par\nobreak\medskip
\begin{center}\fontsize{10}{12}\selectfont\setlength{\tabcolsep}{7pt}\renewcommand{\arraystretch}{1.15}\begin{tabular}{l|c|c}
 & \multicolumn{2}{c}{Unterschied}\\
1813 & in AR. & in Declin.\\\hline
Juni 28 & $+23''6$ & $+22''4$\\
Aug. 16 & $+35.6$ & $.\;.\;.\;.$\\
19 & $+28.5$ & $+4.5$\\
26 & $+22.4$ & $-4.3$\\
\end{tabular}\end{center}

\par\indent Ich bemerke dabei, dass die AR. vom 16. Aug. nur auf einer einzigen Vergleichung beruht,
hingegen die Bestimmungen vom 19. und 26. sich auf zahlreiche und vortrefflich harmonirende
Vergleichungen gründen.


